  \documentclass[11pt,a4paper]{article}
  \usepackage[utf8]{inputenc}
  \usepackage[T1]{fontenc}
  \usepackage[turkish]{babel}
  \usepackage{amsmath}
  \usepackage[margin=2.5cm]{geometry}

  \title{RL Policy Kontrati Kisa Ozeti}
  \author{}
  \date{}

  \begin{document}
  \maketitle

  \section*{Amac}
  Bu notun amaci, MuJoCo Playground tarafindaki ONNX policy kontratini kisa ve net sekilde ozetlemek ve mevcut C++ RL backend ile arasindaki farklari gostermektir.

  \section*{MuJoCo Playground Tarafinda Gercek Kontrat}
  Verilen Python koduna gore policy kontrati sunlardir:

  \begin{itemize}
  \item Girdi adi: \texttt{obs}
  \item Cikti adi: \texttt{continuous\_actions}
  \item Girdi boyutu: \texttt{48}
  \item Cikti boyutu: \texttt{12}
  \end{itemize}

  Inference su sekilde yapilir:
  \[
  \texttt{onnx\_input = \{"obs": obs.reshape(1, -1)\}}
  \]
  \[
  \texttt{onnx\_pred = policy.run(["continuous\_actions"], onnx\_input)[0][0]}
  \]

  \section*{Gozlem Vektoru}
  Python kodundaki gozlem sirasi su sekildedir:

  \[
  \texttt{obs = [linvel,\ gyro,\ gravity,\ joint\_angles,\ joint\_velocities,\ last\_action,\ command]}
  \]

  Acik olarak:

  \begin{enumerate}
  \item \texttt{linvel}: \texttt{data.sensor("local\_linvel").data} \hfill (3 boyut)
  \item \texttt{gyro}: \texttt{data.sensor("gyro").data} \hfill (3 boyut)
  \item \texttt{gravity}: \texttt{imu\_xmat.T @ [0, 0, -1]} \hfill (3 boyut)
  \item \texttt{joint\_angles}: \texttt{data.qpos[7:] - default\_angles} \hfill (12 boyut)
  \item \texttt{joint\_velocities}: \texttt{data.qvel[6:]} \hfill (12 boyut)
  \item \texttt{last\_action}: bir onceki ham policy cikisi \hfill (12 boyut)
  \item \texttt{command}: joystick komutu \hfill (3 boyut)
  \end{enumerate}

  Toplam:
  \[
  3 + 3 + 3 + 12 + 12 + 12 + 3 = 48
  \]

\section*{Action Anlami}
Python kodundaki kontrol su sekildedir:

  \[
  \texttt{data.ctrl[:] = onnx\_pred * action\_scale + default\_angles}
  \]

  Kodda:
  \[
  \texttt{action\_scale = 0.5}
  \]

  Dolayisiyla action yorumu:

  \begin{itemize}
\item Policy cikisi 12 eklem icin \emph{offset} uretir.
\item Hedef eklem acisi:
\[
q_{des} = q_{default} + 0.5 \cdot action
\]
\item Yani bu policy dogrudan eklem pozisyon hedefi uretiyor gibi kullaniliyor.
\end{itemize}

\section*{Default Joint State}
Bu policy icin cok onemli bir nokta, action'in sifir referansi olan \texttt{defaultJointState} bilgisidir. Gozlem tarafinda eklem acilari su sekilde normalize edilir:
\[
q_{rel} = q - q_{default}
\]

Ayni sekilde action da bu poz etrafinda uygulanir:
\[
q_{des} = q_{default} + 0.5 \cdot action
\]

Yani \texttt{defaultJointState} hem observation tarafinda hem de action tarafinda merkez referans pozdur.

Mevcut repodaki \texttt{defaultJointState} sirasi ve degerleri [reference.info] dosyasinda tanimlidir:

\begin{itemize}
\item \texttt{LF\_HAA = -0.10}
\item \texttt{LF\_HFE = -0.70}
\item \texttt{LF\_KFE = 1.30}
\item \texttt{LH\_HAA = 0.10}
\item \texttt{LH\_HFE = -0.30}
\item \texttt{LH\_KFE = 1.20}
\item \texttt{RF\_HAA = 0.10}
\item \texttt{RF\_HFE = 0.70}
\item \texttt{RF\_KFE = -1.30}
\item \texttt{RH\_HAA = -0.10}
\item \texttt{RH\_HFE = 0.30}
\item \texttt{RH\_KFE = -1.20}
\end{itemize}

Bu nedenle yeni bir policy kullanilacaksa sadece \texttt{onnx} dosyasinin degil, policy'nin hangi \texttt{defaultJointState} ile egitildiginin de ayni oldugu teyit edilmelidir.

\section*{Su Anda C++ Tarafinda Dogru Olanlar}
Mevcut C++ RL backend'te su noktalar Python koduyla uyumludur:

  \begin{itemize}
  \item Girdi adi: \texttt{obs}
  \item Cikti adi: \texttt{continuous\_actions}
  \item \texttt{obsDim = 48}
  \item \texttt{actDim = 12}
  \item Gozlem yapisinin genel mantigi dogru:
    lineer hiz, acisal hiz, gravity, eklem acisi, eklem hizi, son action, komut
  \item Action yorumu da dogru yone bakiyor:
    eklem hedefi, default pose etrafinda uretiliyor
  \end{itemize}

  \section*{Su Anda C++ Tarafinda Tam Eslesmeyen Yerler}
  Asil dikkat edilmesi gereken kisim burasidir.

  \subsection*{1. Action scale farki}
  MuJoCo Playground kodunda:
  \[
  \texttt{action\_scale = 0.5}
  \]

  Mevcut \texttt{rl.info} dosyasinda ise:
  \[
  \texttt{actionScale = 0.20}
  \]

  Bu nedenle su an C++ tarafta ayni action cikisi gelse bile ekleme daha kucuk hareket uygulanir.

  \subsection*{2. Hiz kaynagi tam ayni degil}
  Playground kodu su verileri kullaniyor:

  \begin{itemize}
  \item \texttt{local\_linvel}
  \item \texttt{gyro}
  \end{itemize}

  Mevcut C++ backend ise simulator state mesajindan su alanlari kullaniyor:

  \begin{itemize}
  \item \texttt{base\_linvel\_values}
  \item \texttt{base\_angvel\_values}
  \end{itemize}

  Burada problem su:

  \begin{itemize}
  \item \texttt{gyro} yerel IMU acisal hiz bilgisidir.
  \item \texttt{local\_linvel} de robotun yerel eksendeki lineer hizidir.
  \item Bizim mevcut backend ise Python kodundaki sensor isimlerini birebir kullanmiyor.
  \end{itemize}

  Yani boyutlar ayni olsa bile, \emph{hangi frame'den gelen hiz bilgisinin} policy'ye verildigi tam esit olmayabilir.

  \subsection*{3. Joint order kesin teyit edilmeli}
  Python tarafinda eklem sirasi su sekilde geliyor:

  \begin{itemize}
  \item \texttt{data.qpos[7:]}
  \item \texttt{data.qvel[6:]}
  \end{itemize}

  Mevcut repoda kullanilan 12'li sira tipik olarak su:

  \begin{itemize}
  \item \texttt{LF\_HAA, LF\_HFE, LF\_KFE}
  \item \texttt{LH\_HAA, LH\_HFE, LH\_KFE}
  \item \texttt{RF\_HAA, RF\_HFE, RF\_KFE}
  \item \texttt{RH\_HAA, RH\_HFE, RH\_KFE}
  \end{itemize}

  Bu sira buyuk ihtimalle dogru olabilir, fakat policy tarafinda egitim sirasinin birebir ayni oldugu teyit edilmelidir. Aksi halde model yuklenir ama davranis yanlis olur.

  \section*{Pratik Sonuc}
  Bu asamada su sonuca varabiliriz:

  \begin{itemize}
  \item Evet, \texttt{input/output name} ve \texttt{48 -> 12} bilgisi dogrudur.
  \item Evet, action'in ``default pose etrafinda joint offset'' mantigi dogrudur.
  \item Fakat mevcut C++ backend henuz MuJoCo Playground kodunu \emph{birebir} kopyalamis durumda degildir.
  \end{itemize}

  Ozellikle duzeltilmesi gerekenler:

  \begin{enumerate}
  \item \texttt{actionScale} gerekirse \texttt{0.5} yapilmali.
  \item Gozlemde \texttt{local\_linvel} ve \texttt{gyro} ile birebir ayni kaynak/frame kullanilmali.
  \item Joint order kesin olarak teyit edilmeli.
  \end{enumerate}

  \section*{Tek Cumlelik Ozet}
  Su anki framework calisiyor, fakat policy'nin gercekten dogru davranmasi icin C++ tarafindaki observation ve action uygulamasi, MuJoCo Playground kodundaki kontrat ile birebir eslestirilmelidir.

  \end{document}
